\section{序論}

ランサムウェア防御の実務では，検知を速くすることと，バックアップから
復旧できることがしばしば別々の課題として扱われる．しかし，侵入後に
取得されたバックアップが汚染されるなら，両者は同じリスクを別の側面から
扱っている．検知に失敗した感染が長く潜伏すると，直近のバックアップは
すでに汚染され，復旧時に必要なのは「どれだけ新しいバックアップを
持っているか」ではなく「どこまで古い清浄な復旧点へ戻れるか」である．
本稿の出発点は，バックアップを障害後の保険としてではなく，検知投資と
並ぶランサムウェア防御の設計変数として扱うべきだ，という問題意識である．

この問題を難しくするのは，ランサムウェア攻撃では侵入から発症までの
潜伏時間（dwell time）が長く，かつ強い重尾性を持つことが報告されている
\cite{mandiant2023,ibm2023,verizon2023,crowdstrike2023} ためである．
指数分布の尾確率 $\Prob(T>t)=e^{-\mu t}$ は急速に減衰するが，Pareto 分布
\[
  \Prob(T_d>t)=\left(\frac{t_m}{t}\right)^\alpha,\qquad t\ge t_m,
\]
はべき乗則でしか減衰しない．特に $\alpha\le 1$ では $\E[T_d]=\infty$ と
なり，有限平均寿命を暗黙に仮定する古典的な予防保全モデルや
更新報酬モデルをそのまま適用することが難しい
\cite{feller1971,cox1962,barlow1965,nakagawa2005}．実務的には，
「平均的には短いはず」という設計では，まれに長く潜伏した感染が
バックアップ保持期間を越えてから発覚する危険を過小評価する．

既存のセキュリティ投資研究は，投資額とリスク低減の経済的関係を扱う
\cite{gordonloeb2002,hausken2006,laszka2017,fielder2016}．
一方，重尾分布論は極端事象と切断モーメントの扱いに強い
\cite{embrechts1997,resnick2007,bingham1989}．
本稿はこの二つをつなぎ，検知が潜伏時間を指数的に打ち切る効果と，
バックアップ保持期間が全損確率を切り落とす効果を同一モデル内で
比較する．狙いは，尾指数や攻撃頻度を推定すること自体ではなく，
全損時損失，復旧時間，停止損失，保持世代数といった運用上の量から，
バックアップ間隔と清浄バックアップ保持期間にどのような設計指針を
与えられるかを明らかにすることである．

本稿では，早期検知に成功した攻撃の処理費用を
$\Delta$ に依存しない項として分離し，バックアップ間隔に依存する
残余損失を，早期検知に失敗した攻撃の分岐確率から定義する．
この分離により，検知に成功した場合の封じ込め費用と，検知に失敗して
清浄バックアップの有無が問題になる場合の損失を混同しない．
保持世代数を $n$，バックアップ間隔を $\Delta$ とすれば，実効的に
清浄な復旧点へ戻れる期間は $n\Delta$ である．本稿の中心的な問いは，
この復旧可能期間をどの程度確保すべきか，またそのための
バックアップ間隔をどのように決めるべきかである．

主な貢献は以下の通りである．

\begin{enumerate}
  \item 検知レートが正であれば，潜伏状態の平均滞在時間が
  すべての $\alpha>0$ で有限となり，暦時間あたりの費用レートを
  定義できることを示す．
  \item 早期検知確率の低投資域漸近を導き，保持境界なしの残余損失で
  $\alpha=1$ を境界とする相転移が現れる条件を明示する．これは，
  重尾潜伏下でバックアップ保持期間を有限に設計する必要性を示す．
  \item 早期検知失敗分岐の残余期待損失 $\Qfull(\Delta)$ について，
  指数傾き付きの尾確率と切断一次モーメントから大域最小点
  $\Dstar$ を閉形式で与える．この式は，清浄バックアップを残す
  目安期間とバックアップ間隔の運用ルールとして読める．
  \item サービス水準により定まる許容区間上で，固定バックアップ費を
  含む最適間隔 $\Delta^\ast(\lambda)$ が高攻撃頻度極限で
  $\Dstar$ に収束することを証明する．これにより，高頻度攻撃環境で
  尾指数や攻撃頻度に過度に依存しない初期設計基準を与える．
\end{enumerate}

以下では，$\alpha=\alpha(I)$，$\beta=\beta(I)$ と書く場合でも，
投資 $I$ を固定した議論では単に $\alpha,\beta$ と略記する．
確率・期待値は，明記しない限り $T_d$ と $T_{\mathrm{det}}$ の独立性の
もとで評価する．
